%% kapa: exam=1 round=31 year=2020 period=2 form=A subject=appraisal_law subject_ko=감정평가관계법규 seq=4 number=4
%% source: Q-Net 감정평가사 기출문제 (cst00309) artlSeq=5207219 file=extracted/1008/5207219/A형/A형/2020년도 제31회 감정평가사 1차 시험_2교시_A형.pdf
%% conversion: kordoc 4.24.51 -> tools/gen1.py (auto); review=proofread-sonnet; flags=none
\begin{question}{4}
국토의 계획 및 이용에 관한 법령상 용도지역․용도지구․용도구역에 관한 설명으로 옳지 않은 것은?
\begin{choices}
  \item 녹지지역과 공업지역은 도시지역에 속한다.
  \item 용도지구 중 보호지구는 주거 및 교육 환경 보호나 청소년 보호 등의 목적으로 청소년 유해시설 등 특정시설의 입지를 제한할 필요가 있는 지구이다.
  \item 국토교통부장관은 국방부장관의 요청이 있어 보안상 도시의 개발을 제한할 필요가 있다고 인정되면 개발제한구역의 지정을 도시ㆍ군관리계획으로 결정할 수 있다.
  \item 해양수산부장관은 수산자원을 보호ㆍ육성하기 위하여 필요한 공유수면이나 그에 인접한 토지에 대한 수산자원보호구역의 지정을 도시ㆍ군관리계획으로 결정할 수 있다.
  \item 공유수면매립구역이 둘 이상의 용도지역에 걸쳐 있거나 이웃하고 있는 경우 그 매립구역이 속할 용도지역은 도시․군관리계획결정으로 지정하여야 한다.
\end{choices}
\end{question}
\answer{2}
